% ==================== 5.6 本章小结 ====================
\section{本章小结}
\label{sec:ch5_summary}

本章在数字孪生框架下研究了月面岩石识别与避障方法。主要工作和结论如下：

（1）构建了数字孪生平台下的岩石识别框架，包括物理感知层、数据传输层、虚拟处理层和决策应用层四个层次。该框架利用地面计算资源运行复杂的深度学习算法，降低了巡视器的计算负担。

（2）开发了增强型ORB-SLAM2系统，针对月面环境特点进行了光照自适应特征提取、岩石目标检测集成、月面地形约束和月尘干扰抑制等增强设计。该系统能够在月面复杂光照条件下稳定运行，实现岩石的精确识别与定位。

（3）设计了基于岩石识别的避障策略，包括威胁评估、避障决策和路径重规划三个环节。威胁评估综合考虑距离、尺寸、运动方向等因素；避障决策根据威胁度选择停止、绕行或越过策略；路径重规划采用局部修正或全局重规划方法。

（4）通过实验验证了方法的有效性。岩石检测mAP达到94.0\%，定位误差90\%小于0.5 m，SLAM定位精度提升约35\%，避障成功率超过96\%，处理帧率22 fps，满足实时性要求。

本章提出的岩石识别与避障方法为巡视器在复杂月面环境下的安全行走提供了技术支撑。
